
\uline{section donnée de} @@M0@@\uline{), le sous-foncteur de} @@M1@@ \uline{formé des isomorphismes compatibles avec ces rigidifications est représentable par un fermé de} @@M2@@.
\end{enumerate}

En effet, ii) est clair. Quant à i), rappelons-nous que les foncteurs @@M3@@ et @@M4@@ sont représentables par des fibrés vectoriels @@M5@@ et @@M6@@ (d'après EGA III 7.7.8). Alors, si @@M7@@ et @@M8@@, les conditions @@M9@@ et @@M10@@ définissent un fermé @@M11@@ de @@M12@@ et il est clair que @@M13@@.

\textbf{2.5.} Finalement réduisons l'assertion II' à l'assertion suivante:

@@M14@@. \uline{Le théorème} 1.1 \uline{lorsque} @@M15@@ \uline{et} @@M16@@ \uline{sont intègres sur} @@M17@@ \uline{et que la suite}
@@M18@@
\uline{(où} @@M19@@ \uline{est le morphisme structural) est exacte.}

En fait, il suffit pour ceci de prouver qu'il existe une factorisation de @@M20@@
@@M21@@
telle que tout @@M22@@ satisfasse à (*) et, quitte à restreindre @@M23@@, tout @@M24@@ soit intègre sur @@M25@@.

Pour ceci, on pose @@M26@@, @@M27@@ provenant de la factorisation de Stein de @@M28@@ (EGA III 4.3.3), de sorte que @@M29@@, et que @@M30@@ satisfait donc bien à (*). Pour factoriser de même @@M31@@, on utilise le

\textbf{Lemme 2.6.} \uline{Soient} @@M32@@ \uline{un schéma noethérien,} @@M33@@ \uline{un épimorphisme fini (i.e. tel que} @@M34@@ \uline{soit injectif), alors} @@M35@@ \uline{se factorise en une suite finie de morphismes finis} @@M36@@ @@M37@@ \uline{qui sont des épimorphismes effectifs (i.e. tel que la suite}
